\documentclass[11pt]{article}

\usepackage[margin=1in]{geometry}
\usepackage{amsmath,amssymb,amsthm}
\usepackage{array}

\newtheorem{theorem}{Theorem}[section]
\newtheorem{lemma}[theorem]{Lemma}
\newtheorem{proposition}[theorem]{Proposition}
\newtheorem{corollary}[theorem]{Corollary}
\newtheorem{assumption}[theorem]{Assumption}
\theoremstyle{definition}
\newtheorem{definition}[theorem]{Definition}
\theoremstyle{remark}
\newtheorem{remark}[theorem]{Remark}

\newcommand{\area}{\operatorname{area}}
\newcommand{\dinv}{\operatorname{dinv}}
\newcommand{\defc}{\operatorname{defc}}
\newcommand{\UP}{\operatorname{UP}}
\newcommand{\DOWN}{\operatorname{DOWN}}
\newcommand{\East}{\operatorname{East}}
\newcommand{\West}{\operatorname{West}}
\newcommand{\inj}{\operatorname{inj}}
\newcommand{\bk}{\operatorname{bk}}
\newcommand{\Cat}{\operatorname{Cat}}
\newcommand{\NN}{\mathbb N}
\newcommand{\Save}{\mathcal S}
\newcommand{\Repack}{\mathcal R}

\title{The Length-Five Modulo-One Lower-Half Decomposition}
\author{}
\date{}

\begin{document}
\maketitle


\section{Statement and conventions}

Fix an integer
\[
  \tau\ge2,
  \qquad s=5,
  \qquad r=5\tau+1,
  \qquad M=\tau\binom52=10\tau,
  \qquad B=3\tau-1.
\]
A \emph{normalized step-$\tau$ Dyck word} is a word
\[
  D=(x_0,x_1,x_2,x_3,x_4)
\]
of nonnegative integers such that
\[
  x_0=0,
  \qquad x_{i+1}\le x_i+\tau.
\]
Its statistics are
\[
  \area(D)=\sum_i x_i
\]
and
\[
  \dinv_\tau(D)=\sum_{i<j}d_\tau(x_i,x_j),
\]
where
\[
 d_\tau(u,v)=
 \begin{cases}
   (u+\tau-v)_+,&u\le v,\\
   (v+\tau+1-u)_+,&u>v.
 \end{cases}
\]
Here $z_+=\max(0,z)$.  We put
\[
  \defc_\tau(D)=M-\area(D)-\dinv_\tau(D).
\]
For a fixed defect $d$, write
\[
  L_d=\left\lfloor\frac{M-d}{2}\right\rfloor.
\]
Thus $D$ is in the weak lower half of its defect layer precisely when
$\area(D)\le L_d$, and it is strictly below the middle when
$\area(D)\le L_d-1$.

We write $A:B$ for concatenation.  An occurrence $x_j=e>0$ is
\emph{extractable} if exactly one earlier occurrence lies in
\[
  [\max(0,e-\tau),e)
\]
and deleting $x_j$ leaves a normalized word.  If $0<j<4$, the latter
condition is the splice inequality
\[
  x_{j+1}\le x_{j-1}+\tau.
\]
The extraction convention always chooses the leftmost extractable
occurrence.

Conversely, if $e>0$ and a normalized word contains an occurrence in
\[
  [\max(0,e-\tau),e),
\]
then $\inj_e$ inserts $e$ immediately after the first such occurrence.
This is the inverse of extracting the inserted occurrence whenever the
resulting insertion is admissible.

A \emph{full skeleton} is a normalized word with no extractable occurrence.
The two global boundary words
\[
  (0,0,1,0,\tau),
  \qquad
  (0,0,0,0,\tau)
\]
have defect $3\tau=B+1$, so neither occurs in the range considered below.

\subsection{The ordinary local moves}

On a three-window $(a,c,p)$, the map $\East_3$ is defined when
\[
  c\le p+\tau
\]
and is then the identity.  Suppose instead that $c>p+\tau$ in a
five-window $(a,b,c,p,q)$.  Put
\[
  \bk_\tau(b,c)=
  \begin{cases}
    (c,b),&b>c+\tau,\\
    (b,c),&b\le c+\tau.
  \end{cases}
\]
The map $\East_5$ is defined as follows.
\begin{enumerate}
\item If $b\le p+\tau$, it is defined when $b\le q+\tau$ and sends
\[
  (a,b,c,p,q)\longmapsto(a,p,c,b,q).
\]
\item If $b>p+\tau$, write $(b',c')=\bk_\tau(b,c)$.  It is defined when
$c'\le q+\tau$ and sends
\[
  (a,b,c,p,q)\longmapsto(a,p,b',c',q).
\]
\end{enumerate}
The maps $\West_3$ and $\West_5$ are obtained by reversing the input,
applying the corresponding East map, and reversing the output.

The single extraction preceding an East3 test may occur anywhere except the
final position.  There is no stronger position guard at level three.

\begin{assumption}[Full-skeleton lower-half position]
\label{ass:full-lower}
Every full skeleton $S$ with $\defc_\tau(S)\le B$ lies in the weak lower half
of its defect layer:
\[
  \area(S)\le L_{\defc_\tau(S)}.
\]
\end{assumption}

\begin{remark}
The stronger position estimate $\area(S)\le\defc_\tau(S)$ implies
Assumption~\ref{ass:full-lower}.  Indeed, for $d\le B$,
\[
 2\area(S)+d\le3d\le9\tau-3<10\tau=M,
\]
so such a skeleton is in fact strictly below the middle.
\end{remark}

For $0\le d\le B$, let
\[
 \mathcal L_d
 =\{D:\ D\text{ normalized},\ \defc_\tau(D)=d,\ \area(D)\le L_d\}
\]
and let $\mathcal F_d$ be the set of full skeletons of defect $d$.

\begin{theorem}[Length-five lower-half decomposition]
\label{thm:main}
Under Assumption~\ref{ass:full-lower}, there are maps
\[
 \UP:\{D\in\mathcal L_d:\area(D)<L_d\}\longrightarrow\mathcal L_d,
 \qquad
 \DOWN:\mathcal L_d\setminus\mathcal F_d\longrightarrow\mathcal L_d
\]
for every $0\le d\le B$ such that
\[
 \area(\UP(D))=\area(D)+1,
 \qquad
 \dinv_\tau(\UP(D))=\dinv_\tau(D)-1,
\]
\[
 \area(\DOWN(D))=\area(D)-1,
 \qquad
 \dinv_\tau(\DOWN(D))=\dinv_\tau(D)+1,
\]
and the two maps are inverse wherever both sides are defined.  Consequently
\begin{equation}
 \mathcal L_d
 =\bigsqcup_{S\in\mathcal F_d}
   \{\UP^j(S):0\le j\le L_d-\area(S)\}.
 \label{eq:lower-union}
\end{equation}
In particular,
\begin{equation}
 \sum_{D\in\mathcal L_d}q^{\dinv_\tau(D)}t^{\area(D)}
 =\sum_{S\in\mathcal F_d}
   \sum_{j=\area(S)}^{L_d}q^{M-d-j}t^j.
 \label{eq:lower-gf}
\end{equation}
Summing the defect layers gives
\begin{align}
 &\sum_{\substack{D\text{ normalized}\\
                   \defc_\tau(D)\le B\\
                   \area(D)\le L_{\defc_\tau(D)}}}
   q^{\dinv_\tau(D)}t^{\area(D)}\notag\\
 &\qquad=
 \sum_{\substack{S\text{ full skeleton}\\\defc_\tau(S)\le B}}
 \ \sum_{j=\area(S)}^{L_{\defc_\tau(S)}}
 q^{M-\defc_\tau(S)-j}t^j.
 \label{eq:aggregate-lower-gf}
\end{align}
\end{theorem}
\section{Common structural lemmas}

Define the complementary kernel
\[
 K_\tau(u,v)=\tau-d_\tau(u,v)
 =
 \begin{cases}
   \min(\tau,v-u),&u\le v,\\
   \min(\tau,u-v-1),&u>v.
 \end{cases}
\]
For $X=(x_0,\ldots,x_{N-1})$, put
\[
  \lambda_j(X)=\sum_{i<j}K_\tau(x_i,x_j)-x_j.
\]
Then
\begin{equation}
  \defc_\tau(X)=\sum_{j=0}^{N-1}\lambda_j(X),
  \label{eq:defect-lambda}
\end{equation}
where the defect on the left is computed using
$\tau\binom N2$.

\begin{lemma}[Record payment]
\label{lem:record-payment}
Let
\[
  0=r_0<r_1<\cdots<r_h=m,
  \qquad r_i-r_{i-1}\le\tau,
\]
be a chain of record values.  For $0\le u\le m+\tau$,
\[
  \sum_{r_i\le u}K_\tau(r_i,u)\ge u.
\]
Consequently every lambda increment of a normalized word is nonnegative.
If a normalized word begins with two zeros and has record peak $m$, then
\[
  \defc_\tau(X)\ge m.
\]
\end{lemma}

\begin{proof}
Let $r_j$ be the largest record at most $u$.  For $i<j$,
\[
  K_\tau(r_i,u)\ge r_{i+1}-r_i,
\]
while $K_\tau(r_j,u)=u-r_j$.  The sum telescopes to $u$.  In a normalized
word, the preceding records at or below the current entry therefore pay the
subtraction in its lambda increment.

If the word begins with two zeros, the second zero is not needed for these
payments.  At a positive record $r_i$ it contributes
$K_\tau(0,r_i)=\min(\tau,r_i)$.  Thus the unused contributions of the second
zero are at least
\[
  \sum_{i=1}^h\min(\tau,r_i)
  \ge\sum_{i=1}^h(r_i-r_{i-1})=m.
\]
Equation~\eqref{eq:defect-lambda} completes the proof.
\end{proof}

\begin{lemma}[Failed-splice propagation]
\label{lem:propagation}
Let $X=(x_0,\ldots,x_{N-1})$ be normalized and suppose $x_1>0$.  If none
of $x_1,\ldots,x_h$ is extractable, then the entries through $x_{h+1}$ are
strictly increasing records and
\begin{equation}
  x_{i+2}>x_i+\tau
  \qquad(0\le i<h).
  \label{eq:rapid}
\end{equation}
In particular:
\begin{enumerate}
\item every full skeleton of length at least two begins with $(0,0)$;
\item a leftmost extractable occurrence is a record;
\item if appending an entry validly blocks an extractable final occurrence,
the unique-predecessor chain continues to the appended entry, which is
extractable at the new right boundary.
\end{enumerate}
\end{lemma}

\begin{proof}
The occurrence $x_1$ has the unique predecessor $x_0=0$.  If it is not
extractable, its splice fails, so $x_2>x_0+\tau$.  Normalization gives
$x_1\le x_0+\tau$ and $x_2\le x_1+\tau$, whence $x_2>x_1$ and $x_1$ is
the unique predecessor of $x_2$.  The same argument iterates.  At the right
boundary a unique-predecessor occurrence is automatically extractable, which
proves the three consequences.
\end{proof}

\begin{lemma}[Extraction transport]
\label{lem:transport}
Let $e$ be the leftmost extractable occurrence of a normalized word $D$,
let $C$ be obtained by deleting it, and put $p=e-1$.  Then
\begin{align}
 \area(C:p)&=\area(D)-1,\notag\\
 \dinv_\tau(C:p)&=\dinv_\tau(D)+1,\label{eq:transport}\\
 \defc_\tau(C:p)&=\defc_\tau(D).\notag
\end{align}
These identities remain valid even if the final adjacency of $C:p$ fails.
\end{lemma}

\begin{proof}
Write $D=A:e:B$.  For $b\in B$,
\[
  d_\tau(e,b)=d_\tau(b,e-1).
\]
For $a\in A$, the difference
\[
  d_\tau(a,e-1)-d_\tau(a,e)
\]
is $1$ precisely when $a$ lies in the predecessor interval of $e$, is
$-1$ when $e<a\le e+\tau$, and is zero otherwise.  By
Lemma~\ref{lem:propagation}, $e$ is a record, so the negative case does not
occur.  Extractability gives exactly one predecessor and hence a total dinv
increase of $1$.  The area decreases by $1$, proving all three identities.
\end{proof}

\begin{lemma}[Double-zero boundary ledger]
\label{lem:boundary-ledger}
The following four statements hold.
\begin{enumerate}
\item If a normalized length-$s$ word begins with $(0,0)$ and only its final
occurrence is extractable, then its defect is at least $\tau(s-2)$.
\item Let $R$ be a full normalized word of length $s-1$, and let $p$ be the
lowered value arising at a first-extraction stage.  If $R:p$ is invalid or is
not full, then
\[
  \defc_\tau(R:p)\ge\tau(s-2).
\]
\item Let $C:p$ be a first-extraction stage, where $|C|=s-1$, $C$ begins
with $(0,0)$, $C[-1]>p+\tau$, and $C$ has no extractable occurrence outside
its final two positions.  Then
\[
  \defc_\tau(C:p)\ge\tau(s-2).
\]
\item Let $R:p:q$ be a second-extraction stage, where $|R|=s-2$, $R$ begins
with $(0,0)$, $p>R[-1]+\tau$, and $R$ has no extractable occurrence outside
its final position.  Then
\[
  \defc_\tau(R:p:q)\ge\tau(s-2).
\]
\end{enumerate}
\end{lemma}

\begin{proof}
All four assertions are the same kernel ledger.  At every occurrence use the
preceding record chain to pay the negative part of its lambda increment, as
in Lemma~\ref{lem:record-payment}.  The second initial zero is not used in
those payments.  Scan the remaining occurrences from left to right.  If an
occurrence has two or more predecessors, the corresponding two kernel
triangles provide a full $\tau$-unit.  If it has a unique predecessor and is
not extractable, Lemma~\ref{lem:propagation} forces its outgoing splice to
fail.  The unpaid portion of the same $\tau$-unit is carried to the next
record.  Record jumps are at most $\tau$, so these carries are nonnegative
and telescope.

At the right boundary the last carry is closed, respectively, by the
extractable final entry, by the invalid or fullness-obstructing append, by
the far pair $C[-1]>p+\tau$, or by the reverse far pair
$p>R[-1]+\tau$ together with the two extraction anchors.  The ledger contains
$s-2$ full $\tau$-units.  In the two local-window forms the count is
\[
  \underbrace{\tau}_{\text{second zero}}
  +\underbrace{(s-5)\tau}_{\text{interior}}
  +\underbrace{2\tau}_{\text{terminal block}}
  =\tau(s-2).
\]
The first two forms merely move one entry between the interior and terminal
blocks.  All omitted kernel terms are nonnegative.

For the first assertion one may see the count explicitly as follows.  If
$e$ is the final extractable value, then $e\ge\tau+1$.  Pair every nonrecord
middle occurrence $u$ with $0$ and $e$ and use
\[
  K_\tau(0,u)+K_\tau(u,e)\ge\tau.
\]
The record chain supplies the same $\tau$-unit for each record, and the
second zero supplies the final unit.  There are $s-2$ units in total.
\end{proof}

\begin{lemma}[Rapid-chain estimates]
\label{lem:rapid}
Let
\[
  R=(r_0,\ldots,r_{m-1}),
  \qquad r_0=0,
\]
be strictly increasing and normalized, and suppose
\[
  r_{i+2}>r_i+\tau.
\]
If $c=r_{m-1}$, then
\begin{align}
 \dinv_\tau(R)&=\tau(m-1)-c,\label{eq:rapid-dinv}\\
 \sum_i d_\tau(r_i,u)&\le2\tau
 \quad\text{for every }u\ge0.\label{eq:rapid-target}
\end{align}
If $m\ge5$, then
\begin{equation}
  \area(R)>\tau(m-1)\ge\dinv_\tau(R).
  \label{eq:rapid-area}
\end{equation}
If $m\ge3$, then
\begin{equation}
  \area(R)+3c+2>\tau(m+1).
  \label{eq:short-rapid}
\end{equation}
\end{lemma}

\begin{proof}
Only adjacent pairs of $R$ contribute to dinv, and their contributions
 telescope:
\[
 \sum_{i=0}^{m-2}\bigl(\tau-(r_{i+1}-r_i)\bigr)
 =\tau(m-1)-c.
\]
For a fixed target $u$, split the contributing records at $u$.  On each side
their contributions form a truncated triangle of slope one.  The rapid
condition prevents three records on the same side from occupying an interval
of width $\tau$; pairing the two closest possible records on the two sides
gives the bound $2\tau$.

Integrality gives
\[
  r_{2k}\ge k(\tau+1),
  \qquad r_{2k+1}\ge1+k(\tau+1),
\]
which implies~\eqref{eq:rapid-area}.  For~\eqref{eq:short-rapid}, the minimal
length-three chain $(0,1,\tau+1)$ has surplus $7$.  If a chain ending at $c$
is extended by $c'>c$, the surplus in~\eqref{eq:short-rapid} increases by
\[
  4c'-3c-\tau\ge c+4-\tau\ge5.
\]
Induction completes the proof.
\end{proof}

\section{Preliminary UP lemmas}

Throughout this section, let $D=(x_0,\ldots,x_{s-1})$ be normalized and
suppose
\begin{equation}
  0\le d=\defc_\tau(D)\le\tau(s-2)-1,
  \qquad \area(D)\le L_d-1.
  \label{eq:up-hyp}
\end{equation}
Then
\begin{equation}
  \dinv_\tau(D)-\area(D)\ge2.
  \label{eq:strict-lower-gap}
\end{equation}

\begin{lemma}[Validity of the full-skeleton branch]
\label{lem:up-skeleton}
Suppose $D$ is full.  Then $x_{s-1}+1$ has an injection anchor in the prefix.
Consequently the full-skeleton branch of $\UP$ is defined and produces a
normalized word.
\end{lemma}

\begin{proof}
Put $a=x_{s-1}$, let $R=(x_0,\ldots,x_{s-2})$, and let $m=\max R$.
Suppose no prefix entry lies in
\[
  [\max(0,a+1-\tau),a].
\]
Then $a\ge\tau$.  Moreover every prefix entry is at most $a$: the first
entry greater than $a$, if one existed, would have an immediate predecessor
in the forbidden interval.  Hence $m\le a-\tau$.  The final Dyck inequality
now gives
\[
  a\le x_{s-2}+\tau\le m+\tau\le a,
\]
so
\begin{equation}
  a=m+\tau,
  \qquad x_{s-2}=m.
  \label{eq:max-final-jump}
\end{equation}

By Lemma~\ref{lem:propagation}, $D$ begins with two zeros.  The prefix lambda
sum is at least $m$ by Lemma~\ref{lem:record-payment}.  Every prefix entry is
at most $a-\tau$, so every prefix-to-final kernel equals $\tau$.  Therefore
\[
  \lambda_{s-1}(D)=\tau(s-1)-a=\tau(s-2)-m.
\]
It follows that
\[
  \defc_\tau(D)\ge m+\tau(s-2)-m=\tau(s-2),
\]
contrary to~\eqref{eq:up-hyp}.

Thus the anchor exists.  Inserting $a+1$ after the first anchor gives a valid
new left adjacency.  The old Dyck inequality gives the new right adjacency,
and all other adjacencies are unchanged.
\end{proof}

\begin{lemma}[A nonfinal first extraction]
\label{lem:up-first-position}
If $D$ is not full, then it has an extractable occurrence outside its final
position.  In particular its leftmost extractable occurrence is nonfinal.
\end{lemma}

\begin{proof}
Suppose the only extractable occurrence is final.  If $x_1=0$, the first
part of Lemma~\ref{lem:boundary-ledger} gives
\[
  \defc_\tau(D)\ge\tau(s-2),
\]
a contradiction.  If $x_1>0$, Lemma~\ref{lem:propagation} makes the whole
word rapid, and Lemma~\ref{lem:rapid} gives
$\area(D)>\dinv_\tau(D)$, contradicting~\eqref{eq:strict-lower-gap}.
\end{proof}


\section{Length-five East totality and UP classification}

\begin{lemma}[Ordinary East5 is automatic at length five]
\label{lem:east5-total}
Let
\[
 \Sigma=(0,b,c,p,q)
\]
be an East5 stage for which $(0,b,c)$ is normalized and East3 has failed:
\[
 0\le b\le\tau,
 \qquad c\le b+\tau,
 \qquad c>p+\tau.
\]
If $\East_5$ is not defined on $\Sigma$, then
\[
 \defc_\tau(\Sigma)\ge3\tau,
\]
where the defect is computed with the length-five constant $10\tau$.
Consequently every East5 stage arising from a source of defect at most
$B=3\tau-1$ is in the ordinary East5 domain.
\end{lemma}

\begin{proof}
If East5 fails, inspection of its two cases gives one of the following three
possibilities:
\[
\begin{array}{c|c|c}
\text{case}&\text{additional failure}&\text{far pairs}\\ \hline
b\le p+\tau&b>q+\tau&(c,p),(b,q)\\
b>p+\tau,\ b\le c+\tau&c>q+\tau&(b,p),(c,q)\\
b>p+\tau,\ b>c+\tau&b>q+\tau&(c,p),(b,q).
\end{array}
\]
Let
\begin{align*}
 U={}&K_\tau(b,c)+K_\tau(b,p)+K_\tau(b,q)\\
    &+K_\tau(c,p)+K_\tau(c,q)+K_\tau(p,q).
\end{align*}
Then
\begin{equation}
 U\ge4\tau+(c-b-1)_+.
 \label{eq:e5-kernel}
\end{equation}
Indeed, in the first row put $r=b-p$.  The inequalities imply
$1\le r\le\tau$, $c>b$, and $q\le p-1$.  The terms belonging to
$(c,p),(b,q),(c,q)$ equal $\tau$, while
\[
 K_\tau(b,c)=c-b,
 \qquad K_\tau(b,p)=r-1,
 \qquad K_\tau(p,q)\ge\tau-r.
\]
In the second row the three terms for $(b,p),(c,p),(c,q)$ equal $\tau$.
The far pair $(b,p)$ gives
\[
 K_\tau(b,q)+K_\tau(p,q)\ge\tau:
\]
if $q\le p$ the first term is already $\tau$; if $p<q<b$, the two
untruncated terms sum to $b-p-1\ge\tau$; and if $q\ge b$, the second term
is $\tau$.  Also
\[
 K_\tau(b,c)\ge(c-b-1)_+.
\]
In the third row the four terms for
$(b,c),(b,p),(b,q),(c,p)$ all equal $\tau$, and the positive part in
\eqref{eq:e5-kernel} is zero.

The same far-pair list gives
\[
 p,q\le\tau-1.
\]
For $p$ this follows from
$p\le c-\tau-1\le b-1$.  The upper member paired far from $q$ is either
$b\le\tau$ or $c\le b+\tau$, and gives the same conclusion; a negative
upper bound merely says that the corresponding subcase is empty.  Hence
\begin{equation}
 K_\tau(0,b)=b,
 \qquad K_\tau(0,p)=p,
 \qquad K_\tau(0,q)=q.
 \label{eq:e5-zero-pay}
\end{equation}

It remains to pay the entry $c$.  If $b=0$, normalization gives
$c\le\tau$, so $K_\tau(0,c)=c$.  If $b>0$, the record $b$ and the target
$c\le b+\tau$ satisfy
\[
 K_\tau(0,c)+K_\tau(b,c)\ge c.
\]
Since $K_\tau(b,c)\le\tau$, equation~\eqref{eq:e5-kernel} yields in either
case
\begin{equation}
 K_\tau(0,c)+U\ge c+3\tau.
 \label{eq:e5-c-pay}
\end{equation}
Using the complementary-kernel defect identity together with
\eqref{eq:e5-zero-pay} and~\eqref{eq:e5-c-pay}, we obtain
\begin{align*}
 \defc_\tau(0,b,c,p,q)
 &=\sum_{i<j}K_\tau(w_i,w_j)-(b+c+p+q)\\
 &\ge3\tau.
\end{align*}
Two extraction transports preserve defect, so a source of defect at most
$3\tau-1$ cannot reach a failed East5 stage.
\end{proof}

\begin{lemma}[UP position dichotomy after East3 failure]
\label{lem:up-boundary-classification}
Assume~\eqref{eq:up-hyp}.  Let $e_1=p+1$ be the first UP extraction, let
$C$ be its four-entry remainder, and suppose East3 fails on $C:p$.
Then exactly one of the following occurs.
\begin{enumerate}
\item The leftmost extraction of $C$ is at index $1$.  This is the ordinary
East5 branch.
\item The leftmost extraction of $C$ is penultimate.  There are integers
$b,z,P,Q$ such that
\begin{equation}
 C=(0,z+1,P,Q),
 \qquad e_1=b+1,
 \label{eq:up-boundary-shape}
\end{equation}
and
\begin{equation}
 0\le b<z\le\tau-1,
 \qquad
 \tau<P,Q\le z+\tau,
 \qquad
 Q>b+\tau.
 \label{eq:boundary-data}
\end{equation}
In this case the successive extractions from the original word are
$b+1,P,Q$.
\end{enumerate}
\end{lemma}

\begin{proof}
Assume the occurrence at index $1$ of $C$ is not extractable.  If $C[1]=0$,
then $C$ begins with two zeros and has no extraction outside its final two
positions.  East3 failure and the third part of
Lemma~\ref{lem:boundary-ledger}, followed by transport, give
$\defc_\tau(D)\ge3\tau$, a contradiction.  Thus $C[1]>0$.
Lemma~\ref{lem:propagation} shows that the first three entries of $C$ form a
rapid chain.

First suppose the penultimate entry is not extractable.  Its unique-anchor
splice to the final entry must fail, so we may write
\[
 C=(0,a,\tau+u,\tau+a+v),
\]
where
\begin{equation}
 1\le u\le a\le\tau,
 \qquad
 1\le v\le\tau+u-a,
 \qquad
 0\le p<a+v.
 \label{eq:up-affine-range}
\end{equation}
The last inequality is East3 failure.  Put
\[
 f(p)=d_\tau(0,p)+d_\tau(a,p)+(p+1-u)_+.
\]
A direct count in the length-five stage $C:p$ gives
\begin{align}
 \area(C:p)&=2\tau+2a+u+v+p,\notag\\
 \dinv_\tau(C:p)&=2\tau-a-v+f(p),
 \label{eq:up-affine-stats}
\end{align}
and therefore
\begin{align}
 \dinv_\tau(C:p)-\area(C:p)
 &=f(p)-(3a+u+2v+p),
 \label{eq:up-affine-gap}\\
 \area(C:p)+\dinv_\tau(C:p)
 &=4\tau+a+u+p+f(p).
 \label{eq:up-affine-total}
\end{align}
One extraction transport has occurred, so strict lower-halfness gives
\begin{equation}
 f(p)-(3a+u+2v+p)\ge4.
 \label{eq:up-affine-gap4}
\end{equation}

There are three cases.  If $p<a$, then
\[
 f(p)=2\tau+1-a+(p+1-u)_+.
\]
When the positive part is zero, equation~\eqref{eq:up-affine-gap4} is
\[
 4a+u+2v+p\le2\tau-3.
\]
Using $a\ge u$, $v\ge1$, and $p\le u-1$, an elementary integer
optimization gives $u+p\le\tau-4$: otherwise
$5u+p\le2\tau-5$ and $u+p\ge\tau-3$ would force $\tau\le2$, while the
original inequality has no solution for $\tau=2$.  When the positive part
is nonzero, the gap inequality becomes
$2a+u+v\le\tau-1$; since $p\ge u$ and $p<a$, it gives
$2p+1\le\tau-4$.  In both subcases,
\[
 \area(C:p)+\dinv_\tau(C:p)\le7\tau-3.
\]
If $a\le p<\tau$, then
\[
 f(p)=2\tau+a-p+1-u.
\]
Equation~\eqref{eq:up-affine-gap4} gives $2a\le\tau-4$, and again the total
is at most $7\tau-3$.  Finally, if $p\ge\tau$, then
\[
 f(p)\le a+p+1-u,
\]
so the gap in~\eqref{eq:up-affine-gap} is negative.  This case is
impossible.

On the other hand, the defect cutoff requires
\[
 \area(C:p)+\dinv_\tau(C:p)
 =M-\defc_\tau(D)\ge7\tau+1,
\]
a contradiction.  Hence the penultimate entry $P$ of $C$ is extractable.
Write
\[
 C=(0,a,P,Q).
\]
Deleting $P$ gives a normalized word $(0,a,Q)$, so $Q\le a+\tau$.
East3 failure gives $Q>p+\tau\ge\tau$, and therefore $Q$ is then a final
extractable occurrence with unique predecessor $a$.  In particular
$p<a$.  Since reversing the first extraction inserts $p+1$ immediately
after the initial zero, the original word is
\begin{equation}
 D=(0,p+1,a,P,Q).
 \label{eq:up-original-boundary}
\end{equation}

It remains to exclude the endpoint equalities that would spoil the boundary
saving map.  Put
\[
 N=\area(D)+\dinv_\tau(D),
 \qquad
 G=\dinv_\tau(D)-\area(D).
\]
The hypotheses give
\[
 N\ge7\tau+1,
 \qquad G\ge2.
\]
A direct pair count gives the following table; after the first row is
excluded one may use $p\le a-2$, and always $P,Q\ge\tau+1$.
\[
\begin{array}{c|c|c}
\text{endpoint}&N&G\\ \hline
p=a-1&6\tau+2a&4\tau-4a-2P\\
P=a+\tau,\ Q<P&5\tau+p+Q+2&3\tau-4a-p-Q\\
P=a+\tau,\ Q=P&6\tau+a+p+1&2\tau-5a-p-1\\
Q=a+\tau,\ P<a+\tau
 &5\tau+p+P+1+E&3\tau-4a-p-P-1+E,
\end{array}
\]
where
\[
 E=(p+1+\tau-P)_+.
\]
In the first row East3 failure and $Q\le a+\tau$ force
$Q=a+\tau$.  The total and gap bounds then force respectively
$2a\ge\tau+1$ and $2a\le\tau-2$.  In the second row,
$p+Q\le2a+\tau-3$ and $p+Q\ge2\tau-1$ force
$2a\ge\tau+2$ and $4a\le\tau-1$.  In the third row,
$a+p\ge\tau$ forces $2a\ge\tau+2$, while the gap gives
$4a\le\tau-3$.  In the last row, if $E=0$, the total gives
$p+P\ge2\tau$ and hence $2a\ge\tau+3$, while the gap gives
$4a\le\tau-3$.  If $E>0$, substitution gives
\[
 N=6\tau+2p+2,
 \qquad
 G=4\tau-4a-2P,
\]
which force $2a\ge\tau+3$ and $2a\le\tau-2$.  Every endpoint is therefore
impossible.

Set
\[
 b=p,
 \qquad z=a-1.
\]
We have proved
\[
 0\le b<z\le\tau-1,
 \qquad
 \tau<P,Q\le z+\tau.
\]
East3 failure is precisely $Q>b+\tau$.  The penultimate extraction is $P$,
and after deleting it the final extraction is $Q$, proving the asserted
successive extraction order.
\end{proof}

\section{Preliminary DOWN lemmas}

Throughout this section, let $D$ be normalized and suppose
\begin{equation}
  0\le d=\defc_\tau(D)\le\tau(s-2)-1,
  \qquad \area(D)\le L_d.
  \label{eq:down-hyp}
\end{equation}
Assume $D$ is not full.  Let $e_1$ be its leftmost
extractable occurrence and let $C_1$ be the normalized remainder.  Skeleton
and extraction tests are always made on the current prefix alone; lowered
entries already moved to the right are not included.

\begin{lemma}[The DOWN skeleton branch]
\label{lem:down-skeleton}
Under the preceding hypotheses,
\[
  C_1\text{ is full}
  \quad\Longleftrightarrow\quad
  C_1:(e_1-1)\text{ is a normalized full skeleton}.
\]
In particular, if $C_1$ is full, then
\[
  e_1-1\le C_1[-1]+\tau.
\]
\end{lemma}

\begin{proof}
Put $p=e_1-1$ and $P=C_1:p$.  If $P$ is full, then $C_1$ is full.  Indeed,
every nonfinal extraction of $C_1$ persists after appending $p$, while if
the only extraction of $C_1$ is final and the append blocks it,
Lemma~\ref{lem:propagation} makes $p$ extractable at the new right boundary.

Conversely, suppose $C_1$ is full.  It begins with $(0,0)$.  If $P$ were
invalid or not full, the second part of Lemma~\ref{lem:boundary-ledger} and
transport would give
\[
  \defc_\tau(D)=\defc_\tau(P)\ge\tau(s-2),
\]
contrary to~\eqref{eq:down-hyp}.  Hence $P$ is normalized and full, including
the asserted final adjacency.
\end{proof}

\begin{lemma}[Existence of the second DOWN extraction]
\label{lem:down-second}
If $C_1$ is not full, then it has an extractable occurrence.  Thus the second
leftmost extraction is defined on $C_1$, with no position restriction.
\end{lemma}

\begin{proof}
This is the definition of a full skeleton.
\end{proof}

Let $e_2$ be this second extracted value, let $C_2$ be the normalized
remainder, and form
\[
  T_1=C_2:(e_1-1):(e_2-1).
\]

\begin{lemma}[A nonfinal extraction after West3 failure]
\label{lem:down-third-position}
If $\West_3$ fails on the final three entries of $T_1$, then $C_2$ has an
extractable occurrence outside its final position.
\end{lemma}

\begin{proof}
Put
\[
  p=e_1-1,
  \qquad q=e_2-1,
  \qquad c=C_2[-1].
\]
West3 failure is
\begin{equation}
  p>c+\tau.
  \label{eq:west3-fail}
\end{equation}
Suppose $C_2$ has no extraction outside its final position.  Two applications
of transport and the weak lower-half inequality give
\begin{equation}
  \dinv_\tau(T_1)-\area(T_1)
  =\dinv_\tau(D)-\area(D)+4\ge4.
  \label{eq:down-gap}
\end{equation}

If $C_2[1]=0$, the fourth part of Lemma~\ref{lem:boundary-ledger} gives
\[
  \defc_\tau(D)=\defc_\tau(T_1)\ge\tau(s-2),
\]
a contradiction.  Hence $C_2[1]>0$.

Lemma~\ref{lem:propagation} makes
\[
  C_2=R=(r_0,\ldots,r_{m-1}),
  \qquad m=s-2,
  \qquad c=r_{m-1},
\]
a rapid chain.  The unique predecessor of $e_1=p+1$ has value at least
$p+1-\tau>c$.  Since every entry remaining in $C_2$ is at most $c$, this
predecessor must be the second extracted value $e_2$.  Therefore
\begin{equation}
  q=e_2-1\ge p-\tau>c.
  \label{eq:q-above-c}
\end{equation}

No pair from $R$ to $p$ contributes to dinv.  The rapid estimates give
\[
 \dinv_\tau(T_1)
 \le\tau(m-1)-c+2\tau+\tau
 =\tau(m+2)-c.
\]
By integrality in~\eqref{eq:west3-fail} and~\eqref{eq:q-above-c},
\[
 \area(T_1)=\area(R)+p+q
 \ge\area(R)+2c+\tau+2.
\]
Equation~\eqref{eq:short-rapid} now gives
\[
 \area(T_1)-\dinv_\tau(T_1)
 \ge\area(R)+3c+2-\tau(m+1)>0,
\]
contradicting~\eqref{eq:down-gap}.
\end{proof}


\section{Failed West5 stages at length five}

\begin{lemma}[DOWN boundary classification]
\label{lem:down-boundary-classification}
Assume~\eqref{eq:down-hyp}.  Suppose the DOWN construction has failed the
West3 test, made its guaranteed third nonfinal extraction, and reached
\[
 T=(0,b,x,y,z).
\]
If ordinary West5 is not defined on $T$, then
\begin{equation}
 0\le b<z\le\tau-1,
 \qquad
 \tau<x,y\le z+\tau.
 \label{eq:down-boundary-bounds}
\end{equation}
Moreover the original DOWN source has the form
\begin{equation}
 D=(0,z+1,x+1,y+1,b),
 \label{eq:down-source-shape}
\end{equation}
and its first three leftmost extractions are
$x+1,y+1,z+1$.
\end{lemma}

\begin{proof}
After the third extraction the remainder has length two, hence is
$(0,b)$ with $0\le b\le\tau$.  The third extracted occurrence was nonfinal
in a three-entry word and therefore occupied index $1$; its unique
predecessor is the initial zero.  Thus
\[
 0\le z\le\tau-1.
\]
West3 failure, successive extraction, and the three possible West5 failure
modes give
\begin{equation}
 x>b+\tau,
 \qquad x\le y+\tau,
 \qquad \tau<y\le z+\tau.
 \label{eq:down-preliminary}
\end{equation}
It remains to prove $x\le z+\tau$.

Suppose instead that
\[
 y=\tau+v,
 \qquad
 x=\tau+z+u,
\]
where
\[
 1\le v\le z,
 \qquad u\ge1,
 \qquad u\le\tau+v-z.
\]
Put
\[
 h=(b-v)_+,
 \qquad
 k=d_\tau(b,z).
\]
A direct count of the ten pairs of $T$ gives
\begin{align}
 \dinv_\tau(T)&=3\tau-b-z-u+2+h+k,
 \label{eq:down-affine-dinv}\\
 \dinv_\tau(T)-\area(T)
 &=\tau-2b-3z-2u-v+2+h+k,
 \label{eq:down-affine-gap}\\
 \area(T)+\dinv_\tau(T)
 &=5\tau+z+v+2+h+k.
 \label{eq:down-affine-total}
\end{align}
Three extraction transports and weak lower-halfness imply that the gap in
\eqref{eq:down-affine-gap} is at least $6$.

If $b\le z$, then
\[
 k=b+\tau-z.
\]
When $b\le v$, we have $h=0$, and the gap inequality is
\[
 b+4z+2u+v\le2\tau-4.
\]
Since $z\ge\max(b,v)$ and $u\ge1$, the contrary inequality
$b+v\ge\tau-3$ would make the left side at least $3\tau-7$; the remaining
small values $\tau\le3$ admit no solution.  Hence $b+v\le\tau-4$.
When $b>v$, we have $h=b-v$, and the gap inequality gives
$2z+u+v\le\tau-2$, hence
$2b\le\tau-4$.  Substitution into~\eqref{eq:down-affine-total} gives in
both subcases
\[
 \area(T)+\dinv_\tau(T)\le7\tau-2.
\]
If $b>z$, then
\[
 h=b-v,
 \qquad
 k=z+\tau+1-b.
\]
The gap inequality gives
\[
 b+z+u+v\le\tau-2.
\]
Since $b\ge z+1$ and $u,v\ge1$, it follows that
$2z\le\tau-5$, and again
\[
 \area(T)+\dinv_\tau(T)\le7\tau-2.
\]

Transport preserves defect, whereas the cutoff requires
\[
 \area(T)+\dinv_\tau(T)=M-\defc_\tau(D)\ge7\tau+1.
\]
This contradiction proves $x\le z+\tau$.  The inequality
$x>b+\tau$ then gives $b<z$, and~\eqref{eq:down-preliminary} yields all of
\eqref{eq:down-boundary-bounds}.

With these bounds in hand, reverse the extraction history.  The occurrence
$z+1$ has only the initial zero before it.  Injecting $y+1$ and then $x+1$
places both immediately after $z+1$, because $x,y\le z+\tau$; the
successive-extraction inequality determines their displayed order.  Hence
the original word has the form~\eqref{eq:down-source-shape}.  Conversely,
in that word $z+1$ is not extractable because deleting it exposes
$x+1>\tau+1$ after the initial zero.  The occurrence $x+1$ has unique anchor
$z+1$ and a valid splice to $y+1$ because $y\le z+\tau$; after its deletion,
$y+1$ has the same unique anchor and a valid splice to $b$.  Finally
$z+1$ is extractable and leaves $(0,b)$.  The leftmost extraction order is
therefore exactly $x+1,y+1,z+1$.
\end{proof}

\subsection{The formal West extension}

\begin{lemma}[Formal reversal of a failed West5 test]
\label{lem:formal-west}
Let
\[
 T=(a,b,x,y,z)
\]
be a West5 stage satisfying the extraction-history inequalities
\[
 x>b+\tau,
 \qquad x\le y+\tau,
 \qquad y\le z+\tau.
\]
An ordinary West5 failure has exactly one of the following modes:
\[
\begin{array}{c|c|c}
\text{mode}&\text{case test}&\text{failed terminal test}\\ \hline
\mathrm{W1}&y\le b+\tau&y>a+\tau\\
\mathrm{W2}&y>b+\tau,\ y\le x+\tau&x>a+\tau\\
\mathrm{W3}&y>b+\tau,\ y>x+\tau&y>a+\tau.
\end{array}
\]
Define
\begin{equation}
 \widetilde{\West}_5(a,b,x,y,z)=
 \begin{cases}
  (a,y,x,b,z),&\mathrm{W1},\\
  (a,x,y,b,z),&\mathrm{W2},\\
  (a,y,x,b,z),&\mathrm{W3}.
 \end{cases}
 \label{eq:formal-west}
\end{equation}
If $F=\widetilde{\West}_5(T)$, then ordinary East5 is defined on $F$ and
\[
 \East_5(F)=T.
\]
\end{lemma}

\begin{proof}
Reverse $T$, inspect the two cases in the definition of East5, and ignore
only the final domain inequality.  Reversing the resulting output gives the
three rows and~\eqref{eq:formal-west}.  In W1 the case inequality
$y\le b+\tau$ selects East5 case~1 on $F$, while the extraction-history
inequality $y\le z+\tau$ supplies its terminal test.  In W2 the bracket
leaves $(x,y)$ in order, and in W3 it switches $(y,x)$ back to $(x,y)$;
in both cases the same history inequality supplies the terminal test.
Thus East5 is defined and returns $T$ mode by mode.
\end{proof}

For a failed length-five stage the first entry is $a=0$.  Write
\[
 F=\widetilde{\West}_5(T)=(0,P,Q,b,z),
\]
where
\begin{equation}
 (P,Q)=
 \begin{cases}
  (y,x),&\mathrm{W1}\text{ or }\mathrm{W3},\\
  (x,y),&\mathrm{W2}.
 \end{cases}
 \label{eq:PQ-modes}
\end{equation}
The mode conditions always give
\begin{equation}
 Q>b+\tau.
 \label{eq:boundary-east3-fail}
\end{equation}
For stages covered by Lemma~\ref{lem:down-boundary-classification},
\eqref{eq:down-boundary-bounds} also gives
\[
 \tau<P,Q\le z+\tau.
\]
The W3 row is then vacuous, but retaining it makes the formal reversal table
exactly symmetric with the ordinary East5 definition.

\subsection{Saving and repacking}

\begin{definition}[Boundary datum]
\label{def:boundary-datum}
A quadruple $(b,z,P,Q)$ is a \emph{boundary datum} if
\[
 0\le b<z\le\tau-1,
 \qquad
 \tau<P,Q\le z+\tau,
 \qquad
 Q>b+\tau.
\]
For
\[
 F=(0,P,Q,b,z)
\]
define the saving operation
\begin{equation}
 \Save(F)=(0,b+1,z+1,P,Q).
 \label{eq:saving}
\end{equation}
Conversely, for
\[
 U=(0,b+1,z+1,P,Q)
\]
define the boundary repacking operation
\begin{equation}
 \Repack(U)=(0,P,Q,b,z)=F.
 \label{eq:repacking}
\end{equation}
The definition of $\Repack$ includes the three extractions
$b+1,P,Q$ and the replacement of the surviving $z+1$ by $z$; it is not an
ordinary extraction transport.
\end{definition}

\begin{lemma}[Boundary normalization and recognition]
\label{lem:boundary-recognition}
Let $(b,z,P,Q)$ be a boundary datum and put
\[
 U=(0,b+1,z+1,P,Q),
 \qquad F=\Repack(U).
\]
Then $U$ is normalized.  Its first extraction is $b+1$, East3 fails after
that extraction, $z+1$ is not extractable, and the next two extractions are
$P,Q$.  The word $F$ is in the ordinary East5 domain.  If
\[
 T=\East_5(F)=(0,b,x,y,z),
\]
then
\[
 x>b+\tau,
 \qquad x\le y+\tau,
 \qquad \tau<x,y\le z+\tau,
\]
and ordinary West5 fails on $T$ with
\[
 \widetilde{\West}_5(T)=F.
\]
\end{lemma}

\begin{proof}
The first three adjacencies of $U$ follow from
\[
 b+1\le\tau,
 \qquad z+1\le b+1+\tau,
 \qquad P\le z+1+\tau.
\]
Since $P,Q$ both lie in the interval $(\tau,z+\tau]$ of width less than
$\tau$, we also have $Q\le P+\tau$.  Thus $U$ is normalized.

The occurrence $b+1$ is extractable at index $1$.  After deleting it,
East3 tests $Q\le b+\tau$, which fails.  Deleting $z+1$ would expose
$P>\tau$ immediately after the initial zero, so $z+1$ is not extractable.
The occurrence $P$ has unique predecessor $z+1$, and deleting it leaves the
normalized word $(0,z+1,Q)$.  The final occurrence $Q$ then has the same
unique predecessor.  Hence the extraction order is $b+1,P,Q$.

The inequality $Q>b+\tau$ is the East3 failure condition for the five-word
$F=(0,P,Q,b,z)$.  If $P\le b+\tau$, East5 case~1 applies and its terminal
test is $P\le z+\tau$.  If $P>b+\tau$, both entries of
$\bk_\tau(P,Q)$ are at most $z+\tau$, so East5 case~2 applies.  Thus $F$ is
in the East5 domain.

The three East5 cases give
\[
\begin{array}{c|c|c}
\text{case on }F&\East_5(F)&\text{failed West mode}\\ \hline
P\le b+\tau&(0,b,Q,P,z)&\mathrm{W1}\\
P>b+\tau,\ P\le Q+\tau&(0,b,P,Q,z)&\mathrm{W2}\\
P>b+\tau,\ P>Q+\tau&(0,b,Q,P,z)&\mathrm{W3}.
\end{array}
\]
In every row $x>b+\tau$, the entries $x,y$ are the pair $P,Q$, and the
corresponding ordinary West5 terminal test fails because $P>\tau$.
Lemma~\ref{lem:formal-west} then returns $F$.
\end{proof}

\begin{lemma}[Boundary statistic ledger]
\label{lem:boundary-stat}
For a boundary datum, put
\[
 F=(0,P,Q,b,z),
 \qquad U=\Save(F).
\]
Then
\begin{equation}
 \area(U)=\area(F)+2,
 \qquad
 \dinv_\tau(U)=\dinv_\tau(F)-2.
 \label{eq:saving-stat}
\end{equation}
Consequently, if $T$ is a failed West5 stage and
$F=\widetilde{\West}_5(T)$, then the boundary DOWN operation has statistic
change $(-1,+1)$, while the inverse boundary UP operation has change
$(+1,-1)$.
\end{lemma}

\begin{proof}
The area identity is immediate.  Every pair not involving the initial zero
has the same dinv contribution in $F$ and $U$.  Simultaneously shifting
$(b,z)$ to $(b+1,z+1)$ preserves their mutual contribution, and for
$H\in\{P,Q\}$ the boundary inequalities give
\[
 d_\tau(H,b)=d_\tau(b+1,H),
 \qquad
 d_\tau(H,z)=d_\tau(z+1,H).
\]
The two pairs with the initial zero each lose one unit, proving
\eqref{eq:saving-stat}.

Three ordinary extraction transports take a DOWN source to $T$, changing
statistics by $(-3,+3)$.  Lemma~\ref{lem:formal-west} and local preservation
show that $T$ and $F$ have the same statistics, while
$F\mapsto U$ changes them by $(+2,-2)$.  The net DOWN change is therefore
$(-1,+1)$.  Conversely, $U\mapsto F$ changes statistics by $(-2,+2)$,
East5 preserves them, and the three reverse injections change them by
$(+3,-3)$, giving the UP change $(+1,-1)$ without using global
invertibility.
\end{proof}

\section{Definition and totality of the maps}

We first record the ordinary level-three convention in a form that will be
used below.

\begin{lemma}[Level-three reconstruction without a position guard]
\label{lem:level3-reconstruction}
Suppose the first UP extraction is nonfinal and write its normalized
remainder as $C=A:c$, with extracted value $e$ and $p=e-1$.  If
$c\le p+\tau$, then
\[
 \inj_{c+1}\bigl(\inj_e(A)\bigr)
\]
is defined and normalized.  Conversely, the ordinary West3 reconstruction
after two successive DOWN extractions is defined and reverses this
operation.  No condition on the position of $e$ is needed beyond
nonfinality.
\end{lemma}

\begin{proof}
Because $e$ was extracted before the final occurrence $c$, injecting $e$
into $A$ recovers the original word without its last entry.  There is always
an anchor for $c+1$ in that recovered prefix.  If every prefix entry is at
most $c$, then $e$ itself is an anchor because
\[
 c\le e+\tau-1.
\]
If some prefix entry exceeds $c$, take the first such entry; its immediate
predecessor lies in $[c+1-\tau,c]$ by normalization.  Thus inserting $c+1$
after the first anchor is normalized.

For DOWN, write the two extracted values as $x+1,y+1$ and the twice-deleted
word as $C$.  A successful West3 stage is $C:x:y$, with
\[
 x\le C[-1]+\tau.
\]
Injecting $y+1$ into $C:x$ uses the same first anchor that it used in $C$.
If $y+1$ was nonfinal after the first extraction, the displayed West3
inequality supplies the new final adjacency to $x$; if it was final, the
successive-extraction inequality $x\le y+\tau$ supplies that adjacency.
In both cases the old extraction splice is preserved, so the inserted
occurrence is extractable and the inverse East3 test succeeds.
\end{proof}

We shall also use the following local recognition property of the ordinary
five-window pair.

\begin{lemma}[Ordinary East5--West5 branch check]
\label{lem:ordinary-five-recognition}
Suppose an ordinary East5 stage arising after two successive leftmost UP
extractions satisfies
\begin{equation}
 \East_5(a,b,c,x,y)=(a,r,j,k,y).
 \label{eq:east5-symbolic-s5}
\end{equation}
Then
\begin{equation}
 j>r+\tau,
 \qquad
 \West_5(a,r,j,k,y)=(a,b,c,x,y).
 \label{eq:east-west-inverse-s5}
\end{equation}
The raised values $y+1,k+1,j+1$ can be injected from right to left into
$(a,r)$, and are then the first three successive extractions in reverse
order.  Conversely, a successful ordinary West5 stage arising after three
successive DOWN extractions admits the two prescribed reverse injections
and is recognized by the ordinary East5 branch.
\end{lemma}

\begin{proof}
The algebraic recognition is a direct check in the two East5 cases.  In
case~1 the output is $(a,x,c,b,y)$.  In case~2 the pair $(b,c)$ is retained
or switched exactly at the gap $b>c+\tau$.  The original East3 failure gives
$j>r+\tau$ in all three resulting subcases, and reversing the same case
returns the input.

For the reconstruction assertions, retain the inequalities supplied by the
two UP extractions: $(a,b,c)$ is normalized, $x+1$ and $y+1$ have their
unique predecessor anchors in the successive prefixes, and their deletion
splices are valid.  In each East5 case the terminal test is exactly the
anchor inequality for the first raised insertion; the case inequality and
the possible $\bk_\tau$ switch give the second anchor; and the two original
splice inequalities give the final adjacency.  Hence the three displayed
injections are admissible and normalized.  Reading the same inequalities in
reverse proves that the inserted entries are the first three successive
extractions.  The successful West5 statement is the reversed check.
\end{proof}

\subsection{The UP map}

Let $D$ be normalized, have defect at most $B$, and satisfy
$\area(D)<L_{\defc_\tau(D)}$.
\begin{enumerate}
\item If $D=C:a$ is full, define
\begin{equation}
 \UP(D)=\inj_{a+1}(C).
 \label{eq:up-full}
\end{equation}
\item Otherwise extract the leftmost value $e_1=p+1$, leaving $C_1$, and
form
\[
 \Sigma_1=C_1:p.
\]
Write $C_1=A:c$.  If East3 succeeds, define
\begin{equation}
 \UP(D)=\inj_{c+1}\bigl(\inj_{p+1}(A)\bigr).
 \label{eq:up-east3}
\end{equation}
\item Suppose East3 fails.  Let $e_2=q+1$ be the leftmost extraction of
$C_1$.
  \begin{enumerate}
  \item If $e_2$ is at index $1$, delete it to obtain $C_2$ and form
  \[
    \Sigma_2=C_2:p:q.
  \]
  Write
  \[
    \East_5(\Sigma_2)=(u_0,u_1,u_2,u_3,u_4).
  \]
  Define the ordinary East5 branch by
  \begin{equation}
   \UP(D)=
   \inj_{u_2+1}\Bigl(
     \inj_{u_3+1}\bigl(
       \inj_{u_4+1}(u_0,u_1)
     \bigr)
   \Bigr).
   \label{eq:up-east5}
  \end{equation}
  \item Otherwise Lemma~\ref{lem:up-boundary-classification} gives
  \[
    C_1=(0,z+1,P,Q),
    \qquad p=b,
  \]
  with boundary datum $(b,z,P,Q)$.  Extract $P$ and then $Q$, form
  \[
    F=\Repack(0,b+1,z+1,P,Q)=(0,P,Q,b,z),
  \]
  and write
  \[
    \East_5(F)=(0,b,x,y,z).
  \]
  Define the boundary branch by
  \begin{equation}
    \UP(D)=
    \inj_{x+1}\Bigl(
      \inj_{y+1}\bigl(
        \inj_{z+1}(0,b)
      \bigr)
    \Bigr).
    \label{eq:up-boundary}
  \end{equation}
  \end{enumerate}
\end{enumerate}

\subsection{The DOWN map}

Let $D$ be normalized, have defect at most $B$, lie in the weak lower half,
and not be full.
\begin{enumerate}
\item Extract the leftmost value $e_1=p+1$, leaving $C_1$.  If $C_1$ is
full, define
\begin{equation}
 \DOWN(D)=C_1:p.
 \label{eq:down-full}
\end{equation}
\item Otherwise extract the leftmost value $e_2=q+1$ of $C_1$, leaving
$C_2$, and form
\[
 T_1=C_2:p:q.
\]
If West3 succeeds, write
\[
 \West_3(T_1)=(v_0,v_1,v_2,v_3,v_4)
\]
and define
\begin{equation}
 \DOWN(D)=\inj_{v_4+1}(v_0,v_1,v_2,v_3).
 \label{eq:down-west3}
\end{equation}
\item Suppose West3 fails.  Extract the guaranteed nonfinal value
$e_3=z+1$ from $C_2$, leaving $C_3$, and form
\[
 T=C_3:p:q:z.
\]
Since $s=5$, write $C_3=(0,b)$, rename $p=x$ and $q=y$, and write
\[
 T=(0,b,x,y,z).
\]
  \begin{enumerate}
  \item If ordinary West5 succeeds, write
  \[
    \West_5(T)=(v_0,v_1,v_2,v_3,v_4)
  \]
  and define
  \begin{equation}
    \DOWN(D)=
    \inj_{v_3+1}\bigl(
      \inj_{v_4+1}(v_0,v_1,v_2)
    \bigr).
    \label{eq:down-west5}
  \end{equation}
  \item If ordinary West5 fails, form
  \[
    F=\widetilde{\West}_5(T)
  \]
  and define the boundary branch by
  \begin{equation}
    \DOWN(D)=\Save(F).
    \label{eq:down-boundary}
  \end{equation}
  \end{enumerate}
\end{enumerate}

\begin{proposition}[Totality and normalization]
\label{prop:totality}
The map $\UP$ is defined and normalized on every word of defect at most $B$
strictly below the middle.  The map $\DOWN$ is defined and normalized on
every weakly lower-half nonfull word of defect at most $B$.
\end{proposition}

\begin{proof}
For UP, Lemmas~\ref{lem:up-skeleton} and~\ref{lem:up-first-position} give the
full branch or a nonfinal first extraction.  Lemma~\ref{lem:level3-reconstruction}
handles East3 success.  After East3 failure,
Lemma~\ref{lem:up-boundary-classification} gives exactly the ordinary
index-$1$ East5 branch or the penultimate boundary branch.  In the ordinary
case Lemma~\ref{lem:east5-total} proves East5 nonfailure and
Lemma~\ref{lem:ordinary-five-recognition} supplies all injections.  In the
boundary case Lemma~\ref{lem:boundary-recognition} proves that $F$ is in the
East5 domain and that the three injections in~\eqref{eq:up-boundary} are
admissible; explicitly they produce
\[
  (0,z+1,x+1,y+1,b).
\]

For DOWN, Lemmas~\ref{lem:down-skeleton} and~\ref{lem:down-second} give the
first two branches.  Lemma~\ref{lem:level3-reconstruction} handles West3
success.  After West3 failure,
Lemma~\ref{lem:down-third-position} gives a nonfinal third extraction.
Ordinary West5 either succeeds, in which case
Lemma~\ref{lem:ordinary-five-recognition} supplies the two injections, or it
fails.  In the latter case Lemma~\ref{lem:down-boundary-classification}
gives the boundary bounds, Lemma~\ref{lem:formal-west} gives $F$, and
Definition~\ref{def:boundary-datum} and
Lemma~\ref{lem:boundary-recognition} show that $\Save(F)$ is normalized.
Thus no local failure case remains.
\end{proof}

\section{Statistics and invertibility}

\begin{lemma}[Local preservation]
\label{lem:local-preservation}
East3 preserves area and dinv.  On its relevant domain, East5 also preserves
area and dinv.
\end{lemma}

\begin{proof}
East3 is the identity.  East5 only rearranges five entries, so it preserves
area.  Every interaction with an earlier prefix occurrence is unchanged
because the multiset of the five entries is unchanged.  It remains to compare
the six internal pairs among the four moved entries.  Substitution of the
two formulas defining $d_\tau$ in the two East5 cases gives
\[
  \sum_{i<j}d_\tau(w_i,w_j)
  =\sum_{i<j}d_\tau(w'_i,w'_j),
\]
where $w,w'$ are the input and output windows.  In the second case,
$\bk_\tau$ fixes or switches the middle pair according to precisely the gap
at which the two formulas exchange.  This proves local dinv preservation.
\end{proof}


\begin{proposition}[Statistic change]
\label{prop:statistics}
Whenever $\UP$ is defined,
\[
 \area(\UP(D))=\area(D)+1,
 \qquad
 \dinv_\tau(\UP(D))=\dinv_\tau(D)-1.
\]
Whenever $\DOWN$ is defined,
\[
 \area(\DOWN(D))=\area(D)-1,
 \qquad
 \dinv_\tau(\DOWN(D))=\dinv_\tau(D)+1.
\]
Both maps preserve defect.
\end{proposition}

\begin{proof}
For every ordinary branch, if $k$ extractions precede the local East or West
move, repeated use of Lemma~\ref{lem:transport} changes stage statistics by
$(-k,+k)$.  The local move preserves statistics by
Lemma~\ref{lem:local-preservation}.  Each raised injection reverses one
transport and contributes $(+1,-1)$.  The ordinary UP branches make $k+1$
reverse injections and therefore have net change $(+1,-1)$.  The ordinary
DOWN branches make $k-1$ reverse injections and therefore have net change
$(-1,+1)$.  The full branches are the cases $k=0$ for UP and $k=1$ for
DOWN.  The boundary changes are exactly those of
Lemma~\ref{lem:boundary-stat}.  Defect preservation follows from
$M-\area-\dinv_\tau$.
\end{proof}

We use the elementary successive-extraction inequality
\begin{equation}
 e_1\le e_2+\tau.
 \label{eq:successive-extraction-s5}
\end{equation}
Indeed, if the reverse inequality held, deleting $e_1$ would not improve
either the predecessor count or splice of $e_2$, unless the two were
adjacent; adjacency gives the original Dyck inequality and again contradicts
the reverse inequality.

\begin{proposition}[Branchwise invertibility]
\label{prop:inverse}
On the domains in Proposition~\ref{prop:totality},
\[
 \DOWN\circ\UP=\mathrm{id},
 \qquad
 \UP\circ\DOWN=\mathrm{id}.
\]
\end{proposition}

\begin{proof}
We pair the four UP branches with the four DOWN branches.

\smallskip
\noindent\emph{Full skeleton.}
Write a full skeleton as $S=C:a$.  The prefix $C$ is full: an extraction in
$C$ would persist in $S$, or its obstruction would propagate to an
extraction of $a$.  UP injects $a+1$ into $C$.  The inserted occurrence is
the first extractable one.  DOWN extracts it, recognizes the full prefix,
and appends $a$.  Conversely, if DOWN takes its full-prefix branch, it
returns a full word $C:a$ by Lemma~\ref{lem:down-skeleton}, and UP reinserts
$a+1$.

\smallskip
\noindent\emph{East3 and West3.}
Write the East3 stage as $P':b:x$.  The UP output is
\[
 \inj_{b+1}\bigl(\inj_{x+1}(P')\bigr).
\]
DOWN extracts $b+1$ and then $x+1$; its full-prefix branch does not apply,
and the normalized word $P':b$ makes West3 succeed.  West3 is the identity,
and the remaining injection recovers the UP source.

Conversely, suppose DOWN reaches a successful West3 stage $P':b:x:y$.  It
injects $y+1$ before the retained $x$.  UP extracts $y+1$, and its full
branch does not apply.  Inequality~\eqref{eq:successive-extraction-s5} gives
$x\le y+\tau$, exactly the East3 condition.  Lemma~\ref{lem:level3-reconstruction}
then recovers the DOWN source.

\smallskip
\noindent\emph{Ordinary East5 and West5.}
After two UP extractions, write the East5 stage as
\[
 (a,b,c,x,y)
\]
and its output as $(a,r,j,k,y)$.  The raised entries $y+1,k+1,j+1$ are
successively extracted from the UP output.  After the second extraction the
West3 window fails because $j>r+\tau$; after the third,
Lemma~\ref{lem:ordinary-five-recognition} makes ordinary West5 return the
East5 stage, and the two remaining reverse injections recover the original
UP source.

Conversely, after three DOWN extractions let ordinary West5 send
$(a,b,x,y,z)$ to $(a,j,k,r,z)$.  DOWN injects $z+1,r+1$ before the retained
$j,k$.  UP extracts $r+1$; its East3 test fails by the converse recognition
inequality.  It then extracts $z+1$, ordinary East5 returns
$(a,b,x,y,z)$, and the three reverse injections recover the DOWN source.

\smallskip
\noindent\emph{Boundary branch.}
Start with a failed-West5 DOWN source.  By
Lemma~\ref{lem:down-boundary-classification}, DOWN extracts
$x+1,y+1,z+1$ and reaches $T=(0,b,x,y,z)$.  Formal West gives
$F=(0,P,Q,b,z)$ and saving gives
\[
 U=(0,b+1,z+1,P,Q).
\]
By Lemma~\ref{lem:boundary-recognition}, UP extracts $b+1$, East3 fails,
$z+1$ is not extractable, and the next two extractions are $P,Q$.
Repacking returns $F$, East5 returns $T$, and the three injections in
\eqref{eq:up-boundary} recover the original DOWN source.

Conversely, start with a boundary UP word $U$.  Repacking and East5 produce
$T=(0,b,x,y,z)$, and the UP output is
\[
 D'=(0,z+1,x+1,y+1,b).
\]
Lemma~\ref{lem:down-boundary-classification} and its final extraction check
show that DOWN extracts $x+1,y+1,z+1$ in that order.  West3 fails because
$x>b+\tau$.  Ordinary West5 fails, and
Lemma~\ref{lem:boundary-recognition} identifies the same failed mode: formal
West returns $F$, while saving returns $U$.  The boundary branch therefore
cannot be captured by either ordinary five-window branch, and the two
boundary operations are inverse.
\end{proof}

\section{Proof of the lower-half decomposition}

\begin{proof}[Proof of Theorem~\ref{thm:main}]
Fix $0\le d\le B$.  Proposition~\ref{prop:totality} defines UP on every
word in $\mathcal L_d$ of area below $L_d$, and DOWN on every nonfull word
in $\mathcal L_d$.  Proposition~\ref{prop:statistics} shows that the maps
preserve defect and move area by one in opposite directions, and
Proposition~\ref{prop:inverse} shows that they are inverse.

Starting from any $D\in\mathcal L_d$, iterate DOWN while it is defined.
Area decreases at each step, so the process terminates.  It can terminate
only at a full skeleton.  Invertibility makes the terminal skeleton and the
entire descent unique.  Conversely, Assumption~\ref{ass:full-lower} places
every $S\in\mathcal F_d$ in $\mathcal L_d$.  Starting from $S$ and iterating
UP gives exactly one word at each area
\[
 \area(S),\area(S)+1,\ldots,L_d.
\]
No two such strings meet, again by invertibility.  This proves the disjoint
union~\eqref{eq:lower-union}.

For a word of defect $d$ and area $j$, the dinv statistic is $M-d-j$.
Summing over the disjoint strings gives~\eqref{eq:lower-gf}; summing over $0\le d\le B$ gives~\eqref{eq:aggregate-lower-gf}.
\end{proof}

\end{document}
